%=============================================================================
\appendices
%=============================================================================

\section{A Protocol for Measuring $\varepsilon_f$ from Users}
\label{app:protocol}

The most exposed assumption of this paper is the tier annotation
$(0.03, 0.15, 0.30)$, and four experiments reduce the exposure without removing
it. Section~\ref{sec:reliability-empirical} measures an adjacent quantity from
the catalogue; Section~\ref{sec:results-ensemble} shows the conclusion survives
a distribution over the values; Section~\ref{sec:results-corruption} shows how
wrong the tier assignment can be before the correction stops paying; and
Section~\ref{sec:results-em} dispenses with the values entirely, recovering
them from interaction logs.

None of that is the same as measuring $\varepsilon_f$ on people, and the last
of the four is the one whose limits are easiest to overstate. Its logs come
from a simulator whose users err at rates drawn from the annotation, so it
establishes that the parameter is \emph{identifiable} from logs of that size
and that recovering it \emph{suffices}, not that human error rates have this
shape. Only observation of people can establish that. We did not run a human
study, and rather than list that as a limitation we specify the study, so that
the gap is a piece of work someone can do rather than a hole in the argument.

\subsection*{Design}

The estimand is not ``how often is a user wrong'', which is unidentifiable
without ground truth about what the user is thinking. It is the crossover
probability of the channel between the item a user has in mind and the answer
they give: $\varepsilon_f = \Pr[\,\tilde{a}_f \neq a_f(c) \mid c\,]$, where
$a_f(c)$ is the catalogue's value and $\tilde{a}_f$ the user's report. This is
identifiable if the item is known to the experimenter and unknown to be under
test by the participant.

\emph{Task.} A participant is shown a track---title, artist, and a
thirty-second audio excerpt---and asked the same one-versus-rest binary
questions the deployed system would ask, in randomised order. Because the
catalogue value $a_f(c)$ is known, each answer is directly scoreable, and
$\hat\varepsilon_f$ is the per-attribute error rate with a Wilson
interval~\cite{wilson1927}.

\emph{Recall versus perception.} A user of the deployed system answers from
memory, not from an excerpt. The excerpt condition therefore measures a lower
bound on $\varepsilon_f$. We would run both arms: a \emph{perception} arm with
the excerpt available, and a \emph{recall} arm in which participants nominate
tracks they claim to know well and answer without hearing them. The contrast
between the arms is itself the interesting quantity for objective attributes,
where the excerpt makes the answer nearly free but recall does not.

\emph{Sample size.} To separate the annotated tiers we need to distinguish
$0.03$ from $0.15$ and $0.15$ from $0.30$. At $\varepsilon = 0.15$, a Wilson
interval of half-width $0.05$ requires roughly $200$ answers per attribute.
With $20$ attributes and $30$ questions per participant, that is $130$--$140$
participants, or fewer if each answers about several tracks. A pilot at $n=25$
would not resolve individual attributes but would resolve the \emph{tier
contrast}, which is what the method actually consumes: pooling within tiers
gives roughly $1{,}300$ answers per tier at $n=25$, ample to reject tier
equality.

\emph{Analysis.} The per-attribute rate is a proportion with a Wilson interval;
the tier contrast is a mixed-effects logistic regression with a random
intercept per participant and per track, since answers from the same person and
about the same track are not independent. The pre-registered test is the tier
main effect. The secondary analysis regresses $\hat\varepsilon_f$ on (i) the
annotated tier, (ii) the measured catalogue discordance $d_f$ of
\eqref{eq:discordance}, and (iii) the bin count of the attribute, which is the
direct test of the psychophysical mechanism of
Section~\ref{sec:psychophysical}.

\subsection*{What the Study Would and Would Not Settle}

It would settle whether the three-tier ordering is real and roughly how far
apart the tiers are, which is the only property of the annotation that the
method uses: by Proposition~\ref{prop:uniform} only the heterogeneity does any
work, by Section~\ref{sec:results-em} an estimate that is correctly ordered but
badly calibrated performs as well as the truth, and by
Section~\ref{sec:results-corruption} the tier assignment may be wrong for most
attributes before the correction stops paying. The study needs to establish an
ordering, not three decimal places, and it is powered accordingly.

It would not settle the point values for a different catalogue, a different
population, or a different interface, and it should not be read as doing so.
The transportable claim in this paper is conditional and has two parts, which
Section~\ref{sec:results-synthetic} separates. Where answerability is
heterogeneous \emph{and} negatively aligned with apparent informativeness,
classical selection is actively misled and the correction pays most. Where it
is heterogeneous but unaligned, classical selection is merely blind and the
correction still pays, provided the answerable attributes carry enough
information to finish the job. Where they do not, no selection rule can route
around them. Three coarse tiers capture most of the available gain in the first
two cases, and nothing helps in the third.

\section{Reproduction}
\label{app:repro}

Every number in this paper is produced by a script in the released repository
and stored as JSON, and the tables are checked against those files
mechanically rather than transcribed by hand. The determinism requirements of
Section~\ref{sec:system}---single-threaded BLAS, externally supplied random
streams---make the runs bitwise reproducible on a given machine and make each
method's results invariant to which other methods are run alongside it.

Two caveats are worth recording. First, results generated before the
single-thread pin was introduced do not reproduce bitwise; they agree within
sampling error, and where such a run is quoted we say so. Second, the budget
sweep, the $\varepsilon$-ensemble and the cross-domain comparisons are
independent runs at their own sample sizes, so their overlapping cells agree
within sampling error rather than exactly. We report the overlaps rather than
suppressing them.
